%==============================================================================%
% Chapter 12 – Conclusion
%==============================================================================%

\chapter{Conclusion}
\label{ch:conclusion}

\section{Summary of Contributions}

This dissertation has presented a comprehensive formal verification of the
Number-Theoretic Transform in the \haetae{} post-quantum signature scheme.
The verification spans four abstraction layers---from \jasmin{} machine code
to the abstract spectral evaluation formula---and is expressed entirely in the
\easycrypt{} proof assistant.

\paragraph{C1 -- Montgomery correctness.}
We proved that the signed Montgomery reduction in \texttt{reduce.jinc}
correctly computes $a \cdot R^{-1} \bmod q$ within the working range of the
NTT, with explicit bounds on all intermediate values.  The proof required
new signed-arithmetic lemmas for \easycrypt{} that are not present in the
standard library and that are expected to be of independent interest for
future cryptographic verifications.

\paragraph{C2 -- Imperative NTT specification.}
We developed a fully proved imperative NTT specification in \texttt{NTT\_Fq.ec},
covering both the forward and inverse NTT with all loop invariants,
termination witnesses, and coefficient-bound invariants.  The specification is
structurally identical to the C reference implementation and the \jasmin{} source,
making the correspondence proofs straightforward once the arithmetic
infrastructure is in place.

\paragraph{C3 -- Abstract spectral specification.}
We developed the algebraic infrastructure in \texttt{NTTFullAlgebra.ec},
including the primitive root properties of $\zeta_0 = 426$ in $\F_{64513}$,
the orthogonality relation for the negacyclic NTT, bit-reversal identities,
and the butterfly-stage decomposition theorem.  This library distinguishes
the \haetae{} full $256$-point transform from the MLKEM incomplete
$128$-point transform, and provides the formal foundation for the
stage-induction proof.

\paragraph{C4 -- Jasmin refinement.}
We proved the two central refinement lemmas in \texttt{RefJasminNTT.ec}:
\begin{itemize}
  \item \ec{poly\_ntt\_core\_ref}: the \jasmin{} forward NTT is observationally
        equivalent to the imperative specification, with coefficient bounds
        progressing from 16-bit inputs to 24-bit outputs.
  \item \ec{poly\_invntt\_core\_ref}: the \jasmin{} inverse NTT is
        observationally equivalent to the imperative inverse NTT specification,
        preserving 16-bit coefficient bounds in the output.
\end{itemize}

\paragraph{C5 -- End-to-end correctness.}
We assembled all layers in \texttt{NTTEndToEnd.ec}, stating the main theorems
\ec{haetae\_jasmin\_ntt\_correct} and \ec{haetae\_jasmin\_invntt\_correct}.
These theorems directly relate the \jasmin{} exported functions to the abstract
spectral NTT map.

\section{Limitations and Remaining Obligations}

\subsection{Stage Induction Gap}

The one remaining proof obligation is the complete machine-checked proof of the
butterfly-stage induction connecting \ec{NTT\_Fq.NTT.ntt} to
\ec{NTTFullSpec.full\_ntt}.  While the mathematical argument is clear
(\Cref{lem:butterfly_stage}), its formalisation in \easycrypt{} requires
significant index-arithmetic manipulation.  The MLKEM verification did not
provide a reusable template, necessitating new infrastructure.  Completing
this proof is the primary remaining task.

\subsection{No AVX2 Vectorisation}

The current \jasmin{} implementation is scalar (no SIMD).  A production
\haetae{} implementation would use AVX2 intrinsics for a significant
speedup.  Verifying the correctness of an AVX2 NTT in \easycrypt{} would
require extending the proof to handle 256-bit registers and vectorised butterfly
operations.

\subsection{Base Multiplication Not Fully Verified}

The file \texttt{hpoly.jazz} also exports \jasm{poly\_basemul\_jazz}, the
pointwise NTT-domain multiplication.  While the \jasmin{} implementation is
available, its end-to-end correctness proof is not yet complete.  This is
the complement to the NTT correctness: together, the NTT and base
multiplication allow proving that \jasm{poly\_basemul\_jazz} correctly
computes polynomial multiplication in $R_q$.

\subsection{Security Proof Integration}

This dissertation establishes functional correctness only.  Connecting the
NTT implementation to the \emph{security} proof of \haetae{} (i.e., the
reduction from DMSIS) would require additional work at the cryptographic
protocol level.

\section{Future Work}

\paragraph{Completing the stage induction.}
The most pressing future work is completing the butterfly-stage induction in
\texttt{NTTFullAlgebra.ec}.  The key insight needed is a formal lemma relating
the loop counter $k$ to the spectral index $\brev_8(2^s + g)$ at each stage.
We conjecture that a more direct formulation---expressing the partial NTT
directly in terms of the loop variables rather than through the bit-reversal
permutation---would simplify the proof considerably.

\paragraph{Vectorised implementation verification.}
Extending the proof to cover an AVX2 implementation would provide a
performance-verified \haetae{} NTT.  The \jasmin{} language has full support
for AVX2 intrinsics, and the \easycrypt{} extraction handles 256-bit types.
The main challenge is reasoning about the parallelism of 8 butterfly
operations executed simultaneously.

\paragraph{Reusable NTT algebra library.}
The algebraic infrastructure developed in \texttt{NTTFullAlgebra.ec} is
specific to $n = 256$ and $q = 64513$.  A parameterised version, expressing
the NTT algebra in terms of an abstract field $\F_q$ and an abstract
$2n$th root of unity, would be reusable across lattice-based schemes.

\paragraph{Full \haetae{} signing and verification.}
The ultimate goal is a fully verified \haetae{} implementation: key generation,
signing, and verification, all the way from the \jasmin{} source to the
DMSIS security proof.  The NTT verification of this dissertation is the
foundational building block for that larger project.

\paragraph{Comparison with Dilithium and FALCON.}
A comparative study of the NTT verification strategies for Dilithium, FALCON,
and \haetae{} would identify opportunities for sharing algebraic infrastructure
and for developing a unified ``NTT verification toolkit'' for
\easycrypt{}/\jasmin{}.

\section{Closing Remarks}

The formal verification of cryptographic implementations is no longer a
distant aspiration---it is an engineering practice that can be applied to
production code.  This dissertation demonstrates that a full-featured
lattice-based NTT, with all its integer arithmetic nuances, can be formally
verified in \easycrypt{} by careful layering of specifications and proofs.

The main lesson is the value of the three-layer architecture: separating the
word-level arithmetic concerns (Layer 1) from the algorithmic concerns
(Layer 2) from the mathematical concerns (Layer 3) makes each proof layer
tractable.  The interface between layers is clean: word arrays vs.\ integer
arrays (Layers 1--2) and loop invariants vs.\ evaluation formulas (Layers 2--3).

As post-quantum cryptography moves from standardisation to deployment, the
tools and techniques developed here---\jasmin{} for high-assurance
implementations, \easycrypt{} for machine-checked proofs, and the three-layer
verification architecture---provide a clear path to verifying not just NTTs
but complete post-quantum signature and key-encapsulation schemes.
